\scanpage{39}\mspage{38}
% L1–4
Pour a), il suffit de prouver que $T_i$ est fermé dans \deleted{$\unsure{T|X}$} $A$ (les $T_j$, $j\ne i$, l'étant aussi, $T_i$ est aussi ouvert \ldots\ ) Si \deleted{\unsure{$T_i$ est f}} i.e. que
\[
\underbrace{\overline{T_i}}_{\inserted{\substack{\text{``adh'' au}\\\text{sens modéré}\\\text{bien sûr!}}}}\,\cap\, A=T_{i},
\]
% L5–6
C'est immédiat

% L7–12
\uline{NB}\qquad Si $T=T_i$ est un ps.tr., son ``adh.'' modérée n'est autre que le ps.tronçon \inserted{fermé} associé ($=T$ si $T$ réduit à un p\textsuperscript{t}) Il suffit de voir ce qui se passe quand $T$ n'est pas fermé, en un \deleted{$a\in$} \deleted{$\unsure{\partial T}$} \deleted{$\mathrm{Fr}(T)=\dot T$ tel que}
\[
a\in\partial T\setminus T\quad\ldots.
\]

% L13–14, b)
b) $T$ \uline{connexe}\quad \inserted{Si $T=\{\unsure{0|a}\}$ trivial. Sinon $T$ un \deleted{\unsure{ps.}} tronçon.} \uline{NB}\quad $(\mathcal M_{L})_T=\mathcal M_T$,
% L15–17
en regardant $T$ comme un tronçon en lui-$\overline{\text{m}}$. \deleted{Alors} Donc OPS $T=L$. Mais Si $X\subset\deleted{L}T$ modéré,
% L18–19
alors \unsure{:} \deleted{$\unsure{X}$} $X=X_{\Phi}=\bigcup_{i\in\unsure{I|F}}T_i$, alors \deleted{Fr} la frontière modérée est la réunion des \deleted{$\unsure{\partial}$} $T_i^\bullet$,
% L20–22
elle est vide sss les $T_i^\bullet$ sont vides. $\checkmark$Mais $T_i^\bullet=\deleted{\text{vide}}\varnothing\Longrightarrow T_i=L$ donc $X=\varnothing$ ($I=\varnothing$) ou $X=L$ ($I\ne\varnothing$) OK.
